\documentclass[10pt,a4paper]{amsart}
\usepackage{amsmath,amssymb,amsthm,graphicx}
% For commutative diagrams we use Xy-pic
\usepackage[all]{xy}

% We use 2cell for 2-commutative diagrams.
\xyoption{2cell}
% Theorem environments.
%
\theoremstyle{plain}
\newtheorem{theorem}[subsection]{Theorem}
\newtheorem{proposition}[subsection]{Proposition}
\newtheorem{lemma}[subsection]{Lemma}

\theoremstyle{definition}
\newtheorem{definition}[subsection]{Definition}
\newtheorem{example}[subsection]{Example}
\newtheorem{exercise}[subsection]{Exercise}
\newtheorem{situation}[subsection]{Situation}

\theoremstyle{remark}
\newtheorem{remark}[subsection]{Remark}
\newtheorem{remarks}[subsection]{Remarks}

\numberwithin{equation}{subsection}

% Macros
%
\def\lim{\mathop{\mathrm{lim}}\nolimits}
\def\colim{\mathop{\mathrm{colim}}\nolimits}
\def\Spec{\mathop{\mathrm{Spec}}}
\def\Hom{\mathop{\mathrm{Hom}}\nolimits}
\def\Ext{\mathop{\mathrm{Ext}}\nolimits}
\def\SheafHom{\mathop{\mathcal{H}\!\mathit{om}}\nolimits}
\def\SheafExt{\mathop{\mathcal{E}\!\mathit{xt}}\nolimits}
\def\Sch{\mathit{Sch}}
\def\Mor{\mathop{\mathrm{Mor}}\nolimits}
\def\Ob{\mathop{\mathrm{Ob}}\nolimits}
\def\Sh{\mathop{\mathit{Sh}}\nolimits}
\def\NL{\mathop{N\!L}\nolimits}
\def\CH{\mathop{\mathrm{CH}}\nolimits}
\def\proetale{{pro\text{-}\acute{e}tale}}
\def\etale{{\acute{e}tale}}
\def\QCoh{\mathit{QCoh}}
\def\Ker{\mathop{\mathrm{Ker}}}
\def\Im{\mathop{\mathrm{Im}}}
\def\Coker{\mathop{\mathrm{Coker}}}
\def\Coim{\mathop{\mathrm{Coim}}}

% Boxtimes
%
\DeclareMathSymbol{\boxtimes}{\mathbin}{AMSa}{"02}

%
% Macros for moduli stacks/spaces
%
\def\QCohstack{\mathcal{QC}\!\mathit{oh}}
\def\Cohstack{\mathcal{C}\!\mathit{oh}}
\def\Spacesstack{\mathcal{S}\!\mathit{paces}}
\def\Quotfunctor{\mathrm{Quot}}
\def\Hilbfunctor{\mathrm{Hilb}}
\def\Curvesstack{\mathcal{C}\!\mathit{urves}}
\def\Polarizedstack{\mathcal{P}\!\mathit{olarized}}
\def\Complexesstack{\mathcal{C}\!\mathit{omplexes}}
% \Pic is the operator that assigns to X its picard group, usage \Pic(X)
% \Picardstack_{X/B} denotes the Picard stack of X over B
% \Picardfunctor_{X/B} denotes the Picard functor of X over B
\def\Pic{\mathop{\mathrm{Pic}}\nolimits}
\def\Picardstack{\mathcal{P}\!\mathit{ic}}
\def\Picardfunctor{\mathrm{Pic}}
\def\Deformationcategory{\mathcal{D}\!\mathit{ef}}
\newenvironment{reference}{}{}
\newenvironment{slogan}{}{}
\newenvironment{history}{}{}
\usepackage{cleveref}
\crefname{originallemma}{original lemma}{original lemmas}
\Crefname{originallemma}{Original lemma}{Original lemmas}
\makeatletter
\newlabel{lemma-irreducible}{{0C73}{source}{}{sourcepointer-1}{}}
\newlabel{lemma-irreducible@cref}{{[originallemma][0][]0C73}{[1][0][]source}{}{}{}}
\makeatother
\crefname{originallemma}{original lemma}{original lemmas}
\Crefname{originallemma}{Original lemma}{Original lemmas}
\makeatletter
\newlabel{lemma-diagonal-negative}{{0C74}{source}{}{sourcepointer-2}{}}
\newlabel{lemma-diagonal-negative@cref}{{[originallemma][0][]0C74}{[1][0][]source}{}{}{}}
\makeatother
\title{Contracting a minus-one index of a numerical type produces a numerical type of the same genus with one vertex fewer}
\author{The Stacks Project Authors}
\date{stacks-project:models.tex:local-snapshot-20261012; GNU FDL 1.2 or later}
\begin{document}
\maketitle
\noindent\textit{Editorial note.} The counted unit is the Stacks Project Lemma lemma-contract at source lines 780-795 together with its proof at source lines 797-859, a straight adjacent pairing in the chapter file. The chapter is the upstream file models.tex; the pool directory carries the longer name stacks-models and is not used as the chapter name anywhere in this card.\par\medskip
\noindent\textit{Editorial note.} One context span carries the definition of a numerical type at lines 645-662 with its five defining conditions, which the retained proof verifies one by one; the span is the complete definition environment from its opening marker to its closing marker, it is context rather than a counted unit, and it does not overlap the counted interval. It is deliberately not declared as a pointer as well, so that no label is defined twice.\par\medskip
\noindent\textit{Editorial note.} The retained proof names two lemmas of the same chapter, the one excluding the irreducible case and the one on the sign of the diagonal entries; both are declared as non-load-bearing navigation pointers with their stable Stacks tags. Note that this lane does not count the second of them, because its own proof is a bare citation; pointing at its statement from another card is nevertheless valid and its interval is not reproduced.\par\medskip
\noindent\textit{Editorial note.} No citation occurs in the retained spans, so no bibliography entry is supplied.\par\medskip
\noindent\textit{Allowed standard prerequisites (editorial).}\begin{itemize}
\item Numerical types, their five defining conditions and the minus-one index. A numerical type is a tuple of an integer n and integer families of multiplicities, intersection numbers, weights and genera subject to five conditions: positivity of multiplicities, weights and genera; symmetry of the intersection matrix with nonnegative off-diagonal entries; no proper nonempty subset of indices that fails to interact with its complement; vanishing of each weighted row sum; and divisibility of every intersection number by the corresponding weight. An index is a minus-one index when its genus is zero and its self-intersection equals the negative of its weight. The author's own definition and the two lemmas of the same chapter that the proof quotes are the allowed inputs.. The counted lemma is the contraction step on which the whole section on the Picard group of a numerical type rests: it produces a smaller type of the same genus whenever a minus-one index is present.
\end{itemize}\medskip
\ifdefined\hypertarget\hypertarget{sourcepointer-1}{}\fi
\noindent\textit{Source reference.} Original lemma 0C73: focused-seeds/AGP/sources/stacks-snapshot-20261011/tags.txt, source lines 15318--15318. Navigation pointer invoked by the retained proof to exclude the case of a single vertex for the contracted type; the cited statement is an allowed input of the same section and its proof is not reproduced inside the span. The number is the stable Stacks tag of the target, looked up in the shared tag index.\par\medskip
\ifdefined\hypertarget\hypertarget{sourcepointer-2}{}\fi
\noindent\textit{Source reference.} Original lemma 0C74: focused-seeds/AGP/sources/stacks-snapshot-20261011/tags.txt, source lines 15319--15319. Navigation pointer invoked by the retained proof for the sign of the diagonal entries of a type with more than one vertex, used together with the previous pointer; the cited statement is an allowed input and its proof is not reproduced inside the span. The number is the stable Stacks tag of the target, looked up in the shared tag index.\par\medskip
\begin{definition}
\label{definition-type}
A {\it numerical type} $T$ is given by
$$
n, m_i, a_{ij}, w_i, g_i
$$
where $n \geq 1$ is an integer and $m_i$, $a_{ij}$, $w_i$, $g_i$
are integers for $1 \leq i, j \leq n$ subject to the following conditions
\begin{enumerate}
\item $m_i > 0$, $w_i > 0$, $g_i \geq 0$,
\item the matrix $A = (a_{ij})$ is symmetric and $a_{ij} \geq 0$
for $i \not = j$,
\item there is no proper nonempty subset $I \subset \{1, \ldots, n\}$
such that $a_{ij} = 0$ for $i \in I$, $j \not \in I$,
\item for each $i$ we have $\sum_j a_{ij}m_j = 0$, and
\item $w_i | a_{ij}$.
\end{enumerate}
\end{definition}

\begin{lemma}
\label{lemma-contract}
Let $n, m_i, a_{ij}, w_i, g_i$ be a numerical type $T$.
Assume $n$ is a $(-1)$-index. Then there is a numerical
type $T'$ given by $n', m'_i, a'_{ij}, w'_i, g'_i$ with
\begin{enumerate}
\item $n' = n - 1$,
\item $m'_i = m_i$,
\item $a'_{ij} = a_{ij} - a_{in}a_{jn}/a_{nn}$,
\item $w'_i = w_i/2$ if $a_{in}/w_n$ even and $a_{in}/w_i$ odd
and $w'_i = w_i$ else,
\item $g'_i =
\frac{w_i}{w'_i}(g_i - 1) + 1 + \frac{a_{in}^2 - w_na_{in}}{2w'_iw_n}$.
\end{enumerate}
Moreover, we have $g = g'$.
\end{lemma}

\begin{proof}
Observe that $n > 1$ for example by Lemma \ref{lemma-irreducible}
and hence $n' \geq 1$. We check conditions (1) -- (5) of
Definition \ref{definition-type} for $n', m'_i, a'_{ij}, w'_i, g'_i$.

\medskip\noindent
Condition (1) is immediate.

\medskip\noindent
Condition (2). Symmetry of $A' = (a'_{ij})$ is immediate
and since $a_{nn} < 0$ by Lemma \ref{lemma-diagonal-negative}
we see that $a'_{ij} \geq a_{ij} \geq 0$ if $i \not = j$.

\medskip\noindent
Condition (3). Suppose that $I \subset \{1, \ldots, n - 1\}$ such that
$a'_{ii'} = 0$ for $i \in I$ and $i' \in \{1, \ldots, n - 1\} \setminus I$.
Then we see that for each $i \in I$ and $i' \in I'$ we have
$a_{in}a_{i'n} = 0$. Thus either $a_{in} = 0$ for all $i \in I$ and
$I \subset \{1, \ldots, n\}$ is a contradiction for property (3) for $T$,
or $a_{i'n} = 0$ for all $i' \in \{1, \ldots, n - 1\} \setminus I$
and $I \cup \{n\} \subset \{1, \ldots, n\}$ is a contradiction for
property (3) of $T$. Hence (3) holds for $T'$.

\medskip\noindent
Condition (4). We compute
$$
\sum\nolimits_{j = 1}^{n - 1} a'_{ij}m_j =
\sum\nolimits_{j = 1}^{n - 1}
(a_{ij}m_j - \frac{a_{in}a_{jn}m_j}{a_{nn}}) =
- a_{in}m_n - \frac{a_{in}}{a_{nn}}(-a_{nn}m_n) = 0
$$
as desired.

\medskip\noindent
Condition (5). We have to show that $w'_i$ divides $a_{in}a_{jn}/a_{nn}$.
This is clear because $a_{nn} = -w_n$ and $w_n | a_{jn}$ and $w_i | a_{in}$.

\medskip\noindent
To show that $g = g'$ we first write
\begin{align*}
g
& =
1 + \sum\nolimits_{i = 1}^n m_i(w_i(g_i - 1) - \frac{1}{2}a_{ii}) \\
& =
1 + \sum\nolimits_{i = 1}^{n - 1} m_i(w_i(g_i - 1) - \frac{1}{2}a_{ii})
-\frac{1}{2}m_nw_n \\
& =
1 +  \sum\nolimits_{i = 1}^{n - 1}
m_i(w_i(g_i - 1) - \frac{1}{2}a_{ii} - \frac{1}{2}a_{in})
\end{align*}
Comparing with the expression for $g'$ we see that it suffices if
$$
w'_i(g'_i - 1) - \frac{1}{2}a'_{ii} =
w_i(g_i - 1) - \frac{1}{2}a_{in} - \frac{1}{2}a_{ii}
$$
for $i \leq n - 1$. In other words, we have
$$
g'_i = \frac{2w_i(g_i - 1) - a_{in} - a_{ii} + a'_{ii} + 2w'_i}{2w'_i} =
\frac{w_i}{w'_i}(g_i - 1) + 1 + \frac{a_{in}^2 - w_na_{in}}{2w'_iw_n}
$$
It is elementary to check that this is an integer $\geq 0$
if we choose $w'_i$ as in (4).
\end{proof}

\end{document}
